% ===========================================================================
% 2026 FPGA 竞赛创新赛道 AMD Track A (LLM4HLS Agent) 备赛学习手册
% 格式：LaTeX（XeLaTeX / LuaLaTeX），ctexbook 文档类
% 目的：供无 Word 渲染器的 AI 直接阅读源码
% ===========================================================================
\documentclass[12pt,a4paper]{ctexbook}

% ---- 页面与字体 ----
\usepackage{geometry}
\geometry{left=3.0cm,right=2.5cm,top=2.54cm,bottom=2.54cm}
\usepackage{fontspec}
\setmainfont{Times New Roman}
\setCJKmainfont{SimSun}
\setCJKsansfont{SimHei}
\setCJKmonofont{SimSun}

% ---- 表格 ----
\usepackage{booktabs}
\usepackage{longtable}
\usepackage{array}
\newcolumntype{L}[1]{>{\raggedright\arraybackslash}p{#1}}

% ---- 链接 ----
\usepackage[colorlinks=true,linkcolor=blue,urlcolor=cyan]{hyperref}

% ---- 颜色 ----
\usepackage{xcolor}
\definecolor{tipbg}{HTML}{EDF3F5}
\definecolor{warnbg}{HTML}{FBECEA}
\definecolor{notebg}{HTML}{EDF3F5}
\definecolor{accent}{HTML}{1B6B7A}
\definecolor{codebg}{HTML}{F4F8FA}

% ---- 提示框 ----
\usepackage{tcolorbox}
\tcbuselibrary{breakable,skins}
\newtcolorbox{tipbox}{colback=tipbg,colframe=accent,left=4mm,right=4mm,top=2mm,bottom=2mm,breakable}
\newtcolorbox{warnbox}{colback=warnbg,colframe=red!60!black,left=4mm,right=4mm,top=2mm,bottom=2mm,breakable}
\newtcolorbox{notebox}{colback=notebg,colframe=gray!60,left=4mm,right=4mm,top=2mm,bottom=2mm,breakable}

% ---- 代码块 ----
\usepackage{listings}
\lstset{
  basicstyle=\small\ttfamily,
  backgroundcolor=\color{codebg},
  frame=single,
  rulecolor=\color{accent},
  framerule=0.5pt,
  aboveskip=6pt,
  belowskip=6pt,
  breaklines=true,
  showstringspaces=false,
  columns=flexible,
}

% ---- 列表 ----
\usepackage{enumitem}
\setlist{nosep,leftmargin=2em}

% ---- 取消章节编号前缀（用中文） ----
\CTEXsetup[number={\chinese{chapter}}]{chapter}
\CTEXsetup[format={\Large\bfseries\sffamily}]{chapter}

% ---- 标题信息 ----
\title{\sffamily\LARGE 2026 FPGA 竞赛 AMD Track A\\备赛学习手册}
\author{Agent 主 \& HLS 域主（2 人队）}
\date{2026 年 7 月 · v1.0}

\begin{document}

\maketitle

\tableofcontents
\newpage

% ======================================================================
\chapter{比赛要做什么？}

\section{一句话讲清楚这个赛道}

传统 FPGA 竞赛，是你在开发板上手写 Verilog/VHDL，做出一个能跑的硬件。这个赛道不是。

Track A（全名 Budgeted End-to-End LLM4HLS Agent）要你做的事情是：写一个程序（一个 AI agent），里面跑一个大语言模型（LLM，类似 GPT/Claude），让它在“工具调用次数有上限”的约束下，自动地、循环地修好并优化一段 AMD Vitis HLS 的 C/C++ 代码。

打个比方：你不是自己亲手去写那段 HLS 代码，而是在造一个“会自动修 HLS 代码的 AI 工程师”。你交的不是硬件，是这个 AI 工程师本身。

\begin{tipbox}
\textbf{一句话定位}\par
赛道 = 给定一段有问题的 HLS C/C++ 代码 + 测试台 + 规格，你的 agent 要在有限预算内，通过反复调用 csim/cosim/synth 工具、解析反馈、修改代码，最终交出“能跑且 PPA 更优”的 HLS 代码。
\end{tipbox}

\section{关键名词速查表}

下面这些词会反复出现，先建立共同语言。你看不懂某个词，就回这张表。

\begin{longtable}{@{}p{0.16\linewidth}p{0.50\linewidth}p{0.30\linewidth}@{}}
\toprule
\textbf{名词} & \textbf{人话解释} & \textbf{为什么重要} \\
\midrule
\endhead
HLS（高层次综合） & AMD Vitis HLS 工具能把 C/C++ 自动转成 FPGA 硬件。你写 C，加 \texttt{\#pragma} 指挥硬件怎么生成，不用写 Verilog。 & 赛道处理的就是 HLS C 代码，不是 RTL。 \\
agent（智能体） & 以 LLM 当大脑、能循环“看→想→动手→看结果”的程序。 & 这是你要交付的核心产品。 \\
csim（C 仿真） & 跑 C 函数，对照测试台验功能对不对。快。 & 正确性的第一道关。 \\
synth / csynth（综合） & 把 C 编译成 RTL，给你延迟、启动间隔 II、资源用量报告。慢。 & PPA 优化的依据。 \\
cosim（协同仿真） & RTL 和测试台一起跑，验综合后的硬件真的对、时序/握手协议没问题。最慢。 & 正确性的最终关，也是协议类 bug 的来源。 \\
PPA & Performance（性能/延迟/吞吐）、Power（功耗）、Area（资源：LUT/FF/DSP/BRAM）。 & 优化的目标，评分维度之一。 \\
pragma & C 代码里的编译指示，如 \texttt{pipeline/unroll/array\_partition/dataflow}。 & agent 修/优代码的主要抓手。 \\
budget（预算） & csim/cosim/synth 每次调用都花“额度”，次数有限。 & agent 不能暴力试，得聪明诊断、精准改。 \\
\bottomrule
\end{longtable}

\section{你的 Agent 要完成的端到端流程}

拿到题目后，agent 要自动完成下面这条链路，且在预算内收手：

\begin{itemize}
\item 读规格与初始代码：搞清接口契约、数据类型、数值容差、设计约束。
\item 改 C/C++ 代码、加/改 pragma。
\item 调用评估接口（csim/cosim/synth）拿反馈。
\item 解析日志和综合报告，诊断问题在哪。
\item 先修正确性（编译过、csim 过、cosim 过），再优化 PPA。
\item 预算耗尽前终止，交出当前最佳代码。
\end{itemize}

\begin{notebox}
\textbf{关键原则}\par
题目明确写明：correctness 优先于 PPA。先把代码修对，再去优化性能/资源。你的 agent 必须遵守这个顺序。
\end{notebox}

\section{题目会给你什么}

每道题会附带以下产物，agent 要会读懂它们：

\begin{longtable}{@{}p{0.16\linewidth}p{0.50\linewidth}p{0.30\linewidth}@{}}
\toprule
\textbf{产物} & \textbf{内容} & \textbf{你要怎么用} \\
\midrule
\endhead
源文件 & 有问题的基线 C/C++ 代码 & agent 读取并修改它 \\
测试台 & 公开正确性测试 + 构建脚本（如 Makefile） & csim 用的就是它 \\
规格说明 & 接口契约、数据类型、数值容差、设计约束 & agent 据此判断“什么算对” \\
目标约束 & FPGA 平台、HLS 工具版本、时钟目标、可选资源上限 & synth 的硬约束 \\
预算配置 & csim/cosim/synth 的最大调用次数，或统一额度 & agent 的止损线 \\
\bottomrule
\end{longtable}

\section{初始代码可能的几种“病”}

题目给的初始代码不会是好的，至少有下面某一种问题，agent 要会诊断：

\begin{itemize}
\item 功能正确但没优化——能跑，但 PPA 差。
\item HLS 设计编译或综合失败——根本跑不起来。
\item 能编译但 csim/cosim/隐藏功能测试失败——功能有 bug。
\item 结构性问题——死锁、流式数据行为错误、资源严重浪费。
\item 其它与 HLS 编译相关的问题。
\end{itemize}

% ======================================================================
\chapter{我们的目标是什么？}

\section{不是“做出来”，是“打高分”}

这个赛道是能力基准题（benchmark），不是创意展。评分看的是你的 agent 在一批隐藏题目上的表现，不是“你有没有交一个 agent”。所以目标别定成“搭一个能跑的 agent”，而要定成“搭一个在隐藏题集上成功率与 PPA 都靠前的 agent”。

\begin{warnbox}
\textbf{心态校准}\par
“搭出能跑的 agent”不难，一个周末能起骨架；“搭出能打高分的 agent”才是真活，难度全在长尾与诊断质量上。把精力投在后者。
\end{warnbox}

\section{真正的评分维度}

虽然最终评分规则要等比赛方发布，但按赛道描述，agent 的表现大致会落在三个维度上。你做任何设计决策，都要问一句“它提升哪个维度”：

\begin{longtable}{@{}p{0.16\linewidth}p{0.50\linewidth}p{0.30\linewidth}@{}}
\toprule
\textbf{维度} & \textbf{含义} & \textbf{怎么提分} \\
\midrule
\endhead
成功率 & 在隐藏题集上，多少题能修到 correctness 全绿 & 日志解析准 + 知识库覆盖广 + 修法对路 \\
PPA 质量 & 最终代码的延迟/II/资源是否更优 & 读综合报告定位瓶颈 + 用对 pragma 变换 \\
预算效率 & 达成上述结果用了多少 csim/cosim/synth 额度 & 便宜工具先跑、贵的省着用、置信度低就停 \\
\bottomrule
\end{longtable}

\section{战略原则：先正确、后优化}

题目写明 correctness 优先于 PPA。落到 agent 设计上就是两阶段：第一阶段只管把代码修对（编译→csim→cosim 全绿），第二阶段才触发 PPA 优化。不要一上来就想着优化一个还跑不起来的设计。

\begin{notebox}
\textbf{一个要留心的反例}\par
既有研究（COEVO）批评过“先正确再 PPA”的串行做法，主张联合优化。赛道规则要求先正确后优化，但读这类工作能让你在设计评分/奖励策略时不被单一范式锁死。
\end{notebox}

\section{本手册覆盖的里程碑}

这本手册只管到一件事：让你们的第一个 Agent 真正跑起来，并完成一次端到端测试。再往后（PPA 优化打磨、冲分）是后续手册的事。里程碑定义见第五章。

% ======================================================================
\chapter{我们怎么分工？}

\section{两人画像与角色定位}

你们是两人队，画像互补，正好对上这个赛道两大块活：

\begin{itemize}
\item \textbf{你（Agent 主）}：机器人方向，有非常好的编程基础和大模型 Agent 落地经验。负责 agent 全部工程。
\item \textbf{队友（HLS 域主）}：信息工程/射频芯片方向，时间充裕，算法经验有但电路需重学、无 Agent 经验、无 FPGA 经验。负责 HLS 全部领域。
\end{itemize}

\begin{tipbox}
\textbf{给队友的话}\par
射频/信息工程的底子离 HLS 比想象中近——定点（ap\_fixed）、量化、采样、FIR/FFT、吞吐、DSP 资源绑定这些你熟。你缺的只是“工具和数电那半”，不是“信号那半”。从信号处理类 kernel 入手会很有手感。
\end{tipbox}

\section{职责矩阵}

\begin{longtable}{@{}p{0.12\linewidth}p{0.42\linewidth}p{0.42\linewidth}@{}}
\toprule
& \textbf{你 · Agent 主} & \textbf{队友 · HLS 域主} \\
\midrule
\endhead
主责 & agent 全部工程：控制循环、日志解析、prompt、预算调度、知识注入、评测接口对接 & HLS 全部领域：学通 Vitis HLS、bug→修法知识库、造本地题、验证改得对、PPA 报告解读 \\
不碰 & 每种 HLS bug 的修法细节（队友给弹药） & agent 内部循环/解析/prompt 细节（你给接口） \\
关键产出 & 端到端可跑 agent + 评测脚本 + 接口 mock + 预算计数 & bug→修法知识库 + 本地复现题集 + correctness/PPA 验证 \\
强项利用 & Agent 落地经验 → 最开放最值钱的部分由你扛 & 时间充裕 → 长尾知识整理这种苦活由你啃 \\
\bottomrule
\end{longtable}

\section{协作接缝：知识库 schema 契约}

两人队最容易在“接口没对齐”上翻车。你们之间的核心接缝是知识库：队友按固定 schema 产出条目，你写检索器消费它。这个 schema 第一周就要定死，之后两人就能并行。

知识库条目建议 schema（每条一个 JSON 对象）：

\begin{lstlisting}[language=]
{
  id:          "KB-017",
  症状关键词:   ["II not achievable", "pipeline"],
  错误签名:     "Loop .* cannot be pipelined",   // 正则/grep 模式
  根因:         "循环携带依赖 / 存储端口不足",
  修法策略:     "检查数组是否需 array_partition",
  代码示例:     "#pragma HLS array_partition variable=a cyclic ...",
  涉及pragma:   ["pipeline", "array_partition"],
  PPA影响:      "通常降低 II、增加 BRAM/寄存器",
  阶段:         "correctness | ppa"
}
\end{lstlisting}

你的检索器拿当前报错日志去匹配“错误签名”，命中后把 top-k 条目注入 LLM 上下文。schema 早一天定死，两人就早一天并行。

\section{沟通节奏建议}

\begin{itemize}
\item 每天 10 分钟同步：各自昨天做了什么、今天做什么、有没有卡住。
\item 每周一次“对齐会”：跑一次本地评测，看 success rate，决定下周重点。
\item 所有知识库条目、题集、接口 mock 都进同一个 git 仓库，互相看得见。
\end{itemize}

% ======================================================================
\chapter{对方在干嘛？}

这一章是给彼此看的“对方说明书”。你不需要会干对方的活，但要知道对方在干什么、什么时候交付、你什么时候依赖他。

\section{给“你（Agent 主）”看：队友在干什么}

队友在干两件对你至关重要的事：

\begin{itemize}
\item \textbf{建弹药库}：把 HLS 常见失败模式整理成 bug→修法知识库（按第三章 schema）。这是你 agent 的“领域脑子”。没有它，你的 LLM 就是瞎猜烧预算。
\item \textbf{造靶子}：故意把官方 example 改坏，造一批“坏代码 + 已知正确答案”的小题，外加真实日志样例。你的 agent 拿这些来训练/评测。
\end{itemize}

\textbf{你依赖队友的节点}：第一周末要拿到“日志样例包”（队友阶段 0 产出）；第二到三周要拿到知识库初版和 3-5 道小题（队友阶段 1 产出）。这些没到，你的 agent 没靶子打。

\section{给“队友（HLS 域主）”看：你在干什么}

你在干三件对队友至关重要的事：

\begin{itemize}
\item \textbf{搭 agent 骨架}：实现“读代码→调工具→解析日志→让 LLM 出 patch→应用→再跑”的循环。这是消费队友弹药库的机器。
\item \textbf{写日志解析器}：把 Vitis HLS 那一堆噪声日志，抽成干净的“错误类别 + 关键字段”。队友的知识库签名，靠你的解析器来匹配。
\item \textbf{搭评测接口 mock}：在比赛真评估接口还没发布前，做一个假接口让 agent 能跑起来；真接口一发布立刻替换。
\end{itemize}

\textbf{队友依赖你的节点}：第二周你要给出“知识库注入接口”和“日志解析器输出格式”，队友才好照着写知识库签名；同时你要给一个能跑的 mock，队友造的题才能被 agent 跑。

\section{两人交付物如何对接}

\begin{longtable}{@{}p{0.10\linewidth}p{0.42\linewidth}p{0.24\linewidth}p{0.20\linewidth}@{}}
\toprule
\textbf{谁交付} & \textbf{交付物} & \textbf{给谁用} & \textbf{时间} \\
\midrule
\endhead
队友 & 日志样例包（编译错/csim失败/cosim失败各1份） & 你 → 设计日志解析器 & 第一周末 \\
你 & 知识库注入接口 + 日志解析输出格式 & 队友 → 写知识库签名 & 第二周中 \\
队友 & 知识库初版（≥10条）+ 3-5道小题 & 你 → 评测 agent & 第三周末 \\
你 & 可跑的 mock 评估接口 & 队友 → 造的题能被跑 & 第二周末 \\
双方 & 端到端测试题集（5-10道）+ 标准答案 & 共同评测 & 第五周 \\
\bottomrule
\end{longtable}

% ======================================================================
\chapter{最后要完成一个什么样的效果？}

\section{端到端效果描述}

里程碑达成时，你们手里有这样一个东西：给 agent 一道有问题的 HLS 题目（源文件 + 测试台 + 规格 + 预算配置），agent 自动读题、改代码、调 csim/cosim/synth、解析反馈、再改……在预算内把代码修到“编译 + csim + cosim 全绿”，并打印出每步用了多少额度、最后交出可复现脚本与测试报告。

\section{“第一个 Agent 跑通并完成测试”的验收标准}

里程碑的定义（Definition of Done），全部满足才算达成：

\begin{itemize}
\item agent 能在 5-10 道覆盖“编译错 / 功能 bug / cosim 协议错”的题上跑完整流程。
\item 其中至少 60\% 的题 correctness 全绿（编译 + csim + cosim 全过）。
\item agent 有预算计数与止损：超限即优雅终止并交当前最佳代码。
\item 至少 1 次，队友的知识库条目被检索命中并实际生效（可在日志里看到）。
\item 有可复现脚本 + 一段演示（录屏或日志），展示“给坏代码 → agent 自动修到全绿”。
\end{itemize}

\section{一个跑通的样例长什么样}

演示场景举例：你给 agent 一段矩阵乘的 HLS 代码，里面故意缺了 array\_partition，导致 csim 能过但 cosim 因端口冲突失败。agent 的运行日志大致是：

\begin{lstlisting}
[step 1] 调 csim          → PASS
[step 2] 调 cosim         → FAIL  (匹配知识库 KB-017: 端口不足)
[step 3] LLM 提议: 加 #pragma HLS array_partition cyclic
[step 4] 调 cosim         → PASS
[预算] 已用 csim=1 cosim=2 synth=0   剩余额度充足
[交付] correctness 全绿，输出最终代码
\end{lstlisting}

这就算“跑通一道”。能在多道题上稳定复现，就算里程碑达成。

% ======================================================================
\chapter{大体上的地图是怎么样的？}

\section{技术栈全景图}

把整个系统画成一张图，你要随时能在脑子里调出它。每一层都标注了“谁负责”。

\begin{tipbox}
\textbf{你的 Agent（Python）—— 你负责}\par
LLM 大脑（商用 API） ⇄ 控制循环（observe-think-act 反思） ⇄ 预算/调度器（csim/cosim/synth 调用计数与止损）
\end{tipbox}

\vspace{4pt}
\begin{center}\large $\downarrow$ 结构化反馈\end{center}
\vspace{4pt}

\begin{tipbox}
\textbf{日志/报告解析器 —— 你负责}\par
从 Vitis HLS 日志抽出：错误类别、错误行、II、latency、资源用量（LUT/FF/DSP/BRAM）
\end{tipbox}

\vspace{4pt}
\begin{center}\large $\downarrow$ 检索症状\end{center}
\vspace{4pt}

\begin{warnbox}
\textbf{知识库（RAG / 规则）—— 队友负责}\par
bug → 修法条目，按“错误签名”可被检索。内容：编译错 / csim 功能 bug / cosim 协议错 / PPA 杠杆。
\end{warnbox}

\vspace{4pt}
\begin{center}\large $\downarrow$ 调用评估接口\end{center}
\vspace{4pt}

\begin{notebox}
\textbf{比赛评估接口 / 本地 Vitis HLS}\par
csim（C 仿真）→ csynth（综合）→ cosim（协同仿真）。反馈：通过/失败日志 + 综合报告（II / latency / 资源）。
\end{notebox}

\vspace{4pt}
\begin{center}\large $\downarrow$ 最终交付\end{center}
\vspace{4pt}

\begin{tipbox}
\textbf{修好且 PPA 优化的 HLS C/C++ 代码}
\end{tipbox}

\section{数据 / 控制流：从题目到交付}

一次完整运行的流向：题目（代码+规格+预算）→ agent 读入 → 进入控制循环 → [选工具→调用→拿反馈→解析→检索知识库→LLM 出 patch→应用] 反复 → correctness 全绿后进 PPA 阶段 → 预算耗尽或达标则终止 → 交付最终代码 + 运行报告。

\section{知识层在哪里}

知识库不是独立模块，它挂在“日志解析器”和“LLM”之间：解析器把日志变成结构化症状，检索器按症状查知识库，把命中的 bug→修法条目塞进 LLM 的 prompt。这一层是你们俩的接缝，也是这个赛道最大的差异化所在。

\section{角色与地图的对应}

\begin{itemize}
\item 你拥有：Agent 主循环、LLM 大脑、日志解析器、预算调度器、评测接口对接。
\item 队友拥有：知识库内容、本地题集、correctness/PPA 验证。
\item 接缝：知识库 schema 与注入接口（第三章已定）。
\end{itemize}

% ======================================================================
\chapter{实现路径是怎么样的？}

\section{四阶段路线图}

本手册聚焦前两到三个阶段，直到里程碑。第四阶段（PPA 冲分）是后续的事，这里列出只为让你知道方向。

\begin{longtable}{@{}p{0.10\linewidth}p{0.10\linewidth}p{0.38\linewidth}p{0.38\linewidth}@{}}
\toprule
\textbf{阶段} & \textbf{时间} & \textbf{目标} & \textbf{产出} \\
\midrule
\endhead
阶段 0 · 认知+环境 & 第 1 周 & 建立全景、装好环境、跑通官方 example & 日志样例包（队友）、API 通路（你） \\
阶段 1 · 最小正确性 agent & 第 2-3 周 & 搭通 mock 闭环，专攻编译错+csim 功能 bug & 可跑骨架 + 评测脚本 + 知识库初版 \\
阶段 2 · 接真接口+cosim & 第 4-5 周 & 接真工具/真接口，扩到 cosim 协议错，加预算止损 & 里程碑：第一个 Agent 跑通并完成测试 \\
阶段 3 · PPA 优化（后续） & 第 6 周起 & 读综合报告优化 PPA，冲分 & 完整提交版 agent \\
\bottomrule
\end{longtable}

\section{每阶段目标、产出、验收}

\subsection*{阶段 0 · 认知 + 环境}

目标：你建立 agent 范式全景；队友跑通 Vitis HLS 官方 example 并产出真实日志样例。验收：你能一段话讲清赛道；队友本地能跑 csim+csynth+cosim 并交出日志包。

\subsection*{阶段 1 · 最小正确性 agent}

目标：用 mock 接口搭通“读代码→调工具→解析→出 patch→再跑”的闭环，先只攻编译错与 csim 功能 bug，不碰 cosim/PPA。验收：在一道简单编译错题上自动修到 csim 通过；本地评测脚本能跑出 success rate。

\subsection*{阶段 2 · 接真接口 + cosim（里程碑）}

目标：mock 换成真评估接口（或本地 Vitis HLS 命令行），扩到 cosim 协议类错误，加预算计数与止损。验收：见第五章 5.2。

\subsection*{阶段 3 · PPA 优化（后续，不在本手册范围）}

目标：让 agent 读综合报告定位瓶颈，用 pragma 变换优化 PPA，并在预算内取舍。这一阶段需要队友的“PPA 优化速查”弹药，本手册先备着不动。

\section{两条线如何并行汇聚}

你和队友几乎全程并行，只在三个汇合点对齐：

\begin{itemize}
\item 汇合点 1（第一周末）：队友交日志样例包 → 你据此设计解析器。
\item 汇合点 2（第三周末）：队友交知识库初版+小题 → 你跑本地评测。
\item 汇合点 3（第五周末）：双方共做端到端测试题集 → 一起跑里程碑演示。
\end{itemize}

\begin{warnbox}
\textbf{不要串行}\par
别等队友全学完 HLS 你才动手——你阶段 1 用 mock 就能独立推进；队友也别等你 agent 好了才造题——造题不依赖 agent。两人各跑各的线，按汇合点对齐即可。
\end{warnbox}

\section{“最小闭环”定义}

阶段 1 的“最小闭环”是最关键的第一个能跑的东西。它的定义极窄，目的是最快跑通：

\begin{itemize}
\item 输入：一道只有编译错或 csim 功能 bug 的简单题。
\item 工具：只用 mock 的 csim（先不碰 cosim/synth）。
\item 循环：读代码 → 调 mock csim → 解析报错 → LLM 出 patch → 应用 → 再调，直到 csim 过。
\item 不做：不碰 cosim、不碰 PPA、不碰预算止损（先硬编码一个大限额）。
\end{itemize}

这个闭环一旦跑通，后面所有能力都是往这条链上“加料”。

% ======================================================================
\chapter{资料地图：摸着石头过河}

\section{怎么用这份资料清单}

下面三类石头按“必看/推荐/选读”标了优先级。建议：你重点看 B（agent）和 C（LLM4HLS 既有工作）；队友重点看 A（HLS）和 C-1（对口工作）。两人都看 C-4 综述建立全景。链接已尽量验证，个别因反爬/地区限制打不开的，换浏览器或挂梯子即可。

\begin{tipbox}
\textbf{最值钱的一块石头}\par
C-1 里那几篇（Automated C/C++ Repair for HLS、ChatHLS、HLS-Eval、HLSFactory）几乎和赛道逐字对口，是你们最该精读的。先看这几篇再做设计，能少走大量弯路。
\end{tipbox}

\section{A. HLS / Vitis HLS 入门资料}

给零 FPGA 基础的人。队友主线，你也翻 UG1399 的报告章节（为了写解析器）。

{\small
\begin{longtable}{@{}L{0.38\linewidth}L{0.38\linewidth}L{0.12\linewidth}L{0.08\linewidth}@{}}
\toprule
\textbf{资料} & \textbf{看它你能学到什么} & \textbf{优先级} & \textbf{验证} \\
\midrule
\endhead
\href{https://docs.amd.com/r/en-US/ug1399-vitis-hls}{AMD Vitis HLS User Guide (UG1399)} & 官方主文档：csim/csynth/cosim 流程、pragma 语义、报告解读 & 必看 & 门户 \\
\href{https://github.com/Xilinx/Vitis-HLS-Introductory-Examples}{Vitis-HLS-Introductory-Examples} & 官方可综合示例库，按 pragma 分类，带 testbench+Tcl & 必看 & 已验证 \\
\href{https://kastner.ucsd.edu/hlsbook/}{Parallel Programming for FPGAs (Kastner)} & 免费开源 HLS 系统入门书，建立并行/展开/流水线/数据流直觉 & 必看 & 已验证 \\
csynth 报告 / cosim 日志说明 & 综合报告(延迟/II/资源)与 cosim 日志的官方读法。见 UG1399 对应章节 & 必看 & 属UG1399 \\
\href{http://www.openedv.com/}{正点原子 FPGA 教程} & 中文数电/FPGA/Verilog 上手，补底层直觉 & 推荐 & 已验证 \\
\href{http://www.embedfire.com/}{野火 FPGA 教程} & 另一套主流中文 FPGA 入门，任选其一 & 推荐 & 已验证 \\
\href{https://github.com/Xilinx/PYNQ}{PYNQ 项目} & Python+Zynq 降低 FPGA 门槛，理解 overlay/AXI 调硬件核 & 选读 & 已验证 \\
\bottomrule
\end{longtable}
}

\section{B. Agent 工程资料}

给你对齐范式用。你有落地经验，快速过一遍即可，重点看 SWE-agent 和 Anthropic 那篇。

{\small
\begin{longtable}{@{}L{0.38\linewidth}L{0.38\linewidth}L{0.12\linewidth}L{0.08\linewidth}@{}}
\toprule
\textbf{资料} & \textbf{看它你能学到什么} & \textbf{优先级} & \textbf{验证} \\
\midrule
\endhead
\href{https://arxiv.org/abs/2210.03629}{ReAct: Synergizing Reasoning and Acting in LLMs} & 推理+行动交替的 agent 范式，赛道循环的鼻祖 & 必看 & 已验证 \\
\href{https://www.anthropic.com/engineering/building-effective-agents}{Anthropic: Building effective agents} & workflow vs agent、何时该/不该用 agent 的工程判断 & 必看 & 已验证 \\
\href{https://github.com/princeton-nlp/SWE-agent}{SWE-agent + SWE-bench} & 自主改代码过测试的 agent，赛道最对口的软件侧类比 & 必看 & 已验证 \\
\href{https://platform.openai.com/docs/guides/function-calling}{OpenAI Function Calling} & LLM 调外部工具的标准接口范式 & 必看 & 反向爬 \\
\href{https://docs.anthropic.com/en/docs/build-with-claude/tool-use/overview}{Anthropic Tool Use} & Claude 工具循环接线手册（若用 Claude） & 推荐 & 地区受限 \\
\href{https://github.com/langchain-ai/langgraph}{LangChain / LangGraph} & 搭 agent loop/状态机框架。取舍：省轮子 vs 抽象遮蔽预算可控性 & 推荐 & 已验证 \\
\href{https://arxiv.org/abs/2303.11366}{Reflexion} & 把失败经验写回记忆再重试，对应 cosim 失败→总结→重修 & 推荐 & 已验证 \\
\bottomrule
\end{longtable}
}

\section{C. LLM 用于硬件 / HLS 的既有工作（最值钱）}

\subsection*{C-1  直接对口：用 LLM agent 自动修/优化 HLS C 代码}

这几篇最该精读，几乎和赛道逐字对口。

{\small
\begin{longtable}{@{}L{0.38\linewidth}L{0.38\linewidth}L{0.12\linewidth}L{0.08\linewidth}@{}}
\toprule
\textbf{资料} & \textbf{看它你能学到什么} & \textbf{优先级} & \textbf{验证} \\
\midrule
\endhead
\href{https://arxiv.org/abs/2407.03889}{Automated C/C++ Program Repair for HLS via LLMs} & 用 LLM 自动修 HLS C/C++，第一优先精读 & 必看 & 已验证 \\
\href{https://arxiv.org/abs/2507.00642}{ChatHLS} & 多 agent 框架做 HLS 自动调试+directive 调优（有仓库） & 必看 & 已验证 \\
\href{https://arxiv.org/abs/2410.07356}{Optimizing HLS Designs with RAG-LLM} & RAG 辅助 HLS 优化迭代，省调用优化思路 & 必看 & 已验证 \\
\href{https://arxiv.org/abs/2606.17128}{Shift-Left HLS Verification via Knowledge-Augmented LLM Agent} & 知识增强 agent 做 HLS 功能一致性验证 & 必看 & 已验证 \\
\href{https://arxiv.org/abs/2504.21187}{LIFT: LLM Pragma Insertion via GNN SFT} & LLM 自动生成性能关键 pragma，PPA 优化最直接参考 & 必看 & 论文验证 \\
\href{https://arxiv.org/abs/2504.12268}{HLS-Eval} & 评测 LLM 在 HLS 任务上的框架，现成评测脚手架（有仓库） & 必看 & 已验证 \\
\href{https://arxiv.org/abs/2405.00820}{HLSFactory} & 可复现 HLS 数据集/工具框架，本地建库基础设施（有仓库） & 必看 & 已验证 \\
\bottomrule
\end{longtable}
}

\subsection*{C-2  方法论与设计空间}

{\small
\begin{longtable}{@{}L{0.38\linewidth}L{0.38\linewidth}L{0.12\linewidth}L{0.08\linewidth}@{}}
\toprule
\textbf{资料} & \textbf{看它你能学到什么} & \textbf{优先级} & \textbf{验证} \\
\midrule
\endhead
\href{https://arxiv.org/abs/2408.06810}{HLSPilot} & 走 HLS-C 而非 RTL 的方法论与流程设计参考 & 推荐 & 已验证 \\
\href{https://arxiv.org/abs/2505.22086}{iDSE} & LLM 在 directive 空间做 DSE/找 Pareto，PPA 搜索策略参考 & 推荐 & 已验证 \\
\href{https://arxiv.org/abs/2508.03558}{SAGE-HLS} & AST/语法感知引导生成可综合 HLS，降低修复次数 & 推荐 & 已验证 \\
\href{https://arxiv.org/abs/2408.10428}{Are LLMs Any Good for HLS?} & 综述+实验，建立“LLM 在 HLS 行不行”全景 & 推荐 & 已验证 \\
\href{https://arxiv.org/abs/2406.09233}{C2HLSC} & LLM 迭代重写普通 C 成可综合 HLS-C，案例可练手 & 推荐 & 已验证 \\
\bottomrule
\end{longtable}
}

\subsection*{C-3  RTL 侧 agent 与反馈循环（可迁移到 HLS）}

{\small
\begin{longtable}{@{}L{0.38\linewidth}L{0.38\linewidth}L{0.12\linewidth}L{0.08\linewidth}@{}}
\toprule
\textbf{资料} & \textbf{看它你能学到什么} & \textbf{优先级} & \textbf{验证} \\
\midrule
\endhead
\href{https://arxiv.org/abs/2311.04887}{AutoChip} & LLM+反馈迭代生成/修 HDL 的经典，范式可直接迁移（有仓库） & 必看 & 已验证 \\
\href{https://arxiv.org/abs/2411.11856}{Auto-Improving LLM Verilog via EDA Tool Feedback} & 多轮 EDA 工具反馈修代码，与赛道同构 & 必看 & 已验证 \\
\href{https://arxiv.org/abs/2311.16543}{RTLFixer} & RAG+ReAct 让 LLM 自主调试修 Verilog 语法，赛道“修正确性”蓝本 & 必看 & 已验证 \\
\href{https://arxiv.org/abs/2604.15001}{COEVO} & 批评“先正确再 PPA”串行，主张联合优化——策略对照 & 必看 & 已验证 \\
\href{https://arxiv.org/abs/2309.07544}{VerilogEval} & Verilog 生成评测基准，借鉴评测口径（有仓库） & 推荐 & 已验证 \\
\href{https://arxiv.org/abs/2308.10204}{ChatEDA} & LLM 规划+EDA 工具执行架构参考（有仓库） & 推荐 & 已验证 \\
\bottomrule
\end{longtable}
}

\subsection*{C-4  综述（建立全景）}

{\small
\begin{longtable}{@{}L{0.38\linewidth}L{0.38\linewidth}L{0.12\linewidth}L{0.08\linewidth}@{}}
\toprule
\textbf{资料} & \textbf{看它你能学到什么} & \textbf{优先级} & \textbf{验证} \\
\midrule
\endhead
\href{https://arxiv.org/abs/2501.09655}{A Survey of LLMs for EDA} & LLM for EDA 综述，快速建立全领域地图与入口 & 必看 & 已验证 \\
\bottomrule
\end{longtable}
}

\section{阅读顺序建议}

\begin{itemize}
\item 第一周：两人都看 C-4 综述建立全景；你看 ReAct + Anthropic 那篇；队友看 Kastner 前几章 + UG1399 报告章节。
\item 第二周：你精读 SWE-agent + C-1 的 Automated Repair 与 ChatHLS；队友精读 C-1 的 HLS-Eval/HLSFactory（找现成脚手架）。
\item 第三周起：按需翻 C-2/C-3，遇到具体问题再查。
\end{itemize}

% ======================================================================
\chapter{待办清单（管到第一个 Agent 跑通 + 完成测试）}

\section{怎么用这两份清单}

两份清单各自独立推进，按第七章三个汇合点对齐。每项都有“完成标志”，达标才打勾。

\section{【你 · Agent 主】待办清单}

\subsection*{阶段 0 · 认知 + 环境（第 1 周）}

\begin{itemize}
\item[☐] 通读本手册第一、六、七章，建立全景。\hfill {\small\textit{完成标志：能用一段话向队友讲清“赛道要干什么、agent 要做什么循环”}}
\item[☐] 读 ReAct 论文 + Anthropic《Building effective agents》。\hfill {\small\textit{完成标志：能画出 observe→think→act 循环，说清 workflow 与 agent 的区别}}
\item[☐] 读 SWE-agent 仓库 README + 论文。\hfill {\small\textit{完成标志：列出 3 条可迁移到 HLS 的设计}}
\item[☐] 装好开发环境 + 选定商用 LLM API。\hfill {\small\textit{完成标志：跑通一个最小“调 API → 拿结构化 JSON 输出”的脚本}}
\item[☐] 向队友要到“Vitis HLS 日志样例包”（依赖队友阶段 0）。\hfill {\small\textit{完成标志：拿到至少 1 份真实日志，能指出错误行/II/资源字段在哪}}
\end{itemize}

\subsection*{阶段 1 · 最小正确性 agent（第 2-3 周）}

\begin{itemize}
\item[☐] 定义评估接口 mock：模拟 csim/cosim/synth 的输入输出。\hfill {\small\textit{完成标志：mock 能据传入代码返回“编译错/csim 过/cosim 失败”等预设结果}}
\item[☐] 实现日志解析器：从 mock 日志抽出错误类型与关键字段。\hfill {\small\textit{完成标志：对 3 类典型日志能正确抽出错误类别}}
\item[☐] 实现 ReAct 主循环：读代码→调 mock→解析→LLM 出 patch→应用→再调。\hfill {\small\textit{完成标志：在一道简单编译错题上自动修到 csim 通过}}
\item[☐] 接入队友的 bug→修法知识库（哪怕先 3 条）：按症状检索注入 prompt。\hfill {\small\textit{完成标志：知识库条目能被检索并出现在 LLM 上下文里}}
\item[☐] 写本地评测脚本：在队友造的 3-5 道小题上跑 success rate。\hfill {\small\textit{完成标志：输出每题“通过/未通过/用了几次工具调用”}}
\end{itemize}

\subsection*{阶段 2 · 接真接口 + cosim（第 4-5 周）→ 里程碑}

\begin{itemize}
\item[☐] 把 mock 替换为比赛真评估接口（规则一发布立刻接）。\hfill {\small\textit{完成标志：agent 能调真 csim/cosim/synth 拿真反馈}}
\item[☐] 或本地接 Vitis HLS 命令行（vitis\_hls -f run.tcl）跑通真 csim/synth。\hfill {\small\textit{完成标志：本地能跑通 1 个真例子的 csim+synth，解析真报告}}
\item[☐] 扩到 cosim 协议类错误：dataflow 死锁 / AXI-Stream 握手错。\hfill {\small\textit{完成标志：在一道 cosim 失败题上自动修到 cosim 通过}}
\item[☐] 加预算计数与止损：记录已用 csim/cosim/synth 次数，超限即停。\hfill {\small\textit{完成标志：agent 在预算耗尽时优雅终止并交当前最佳代码}}
\item[☐] 端到端测试：在 5-10 道覆盖“编译错/功能 bug/cosim 协议错”的题上跑完整流程。\hfill {\small\textit{完成标志：≥60\% 题目 correctness 全绿，并输出测试报告}}
\item[☐] 里程碑：录一段演示——给一道坏 HLS，agent 自动修到 csim+cosim 通过。\hfill {\small\textit{完成标志：有可复现脚本 + 演示录屏/日志}}
\end{itemize}

\section{【队友 · HLS 域主】待办清单}

\subsection*{阶段 0 · 认知 + 环境（第 1 周）}

\begin{itemize}
\item[☐] 刷新数电：组合/时序逻辑、触发器、FSM、时序（建立/保持）、流水线。\hfill {\small\textit{完成标志：能讲清“流水线为什么提高吞吐、寄存器干啥”}}
\item[☐] 读 Kastner 前几章 + UG1399 的 csim/csynth/cosim 章节。\hfill {\small\textit{完成标志：能用一段话讲清 csim/synth/cosim 各验什么、报告里 II/latency/资源在哪}}
\item[☐] 装好 Vitis HLS（或用学校/云端）。\hfill {\small\textit{完成标志：本地跑通 1 个官方 example 的 csim+csynth+cosim}}
\item[☐] 产出“日志样例包”给 Agent 主：编译错/csim 失败/cosim 失败/synth 成功各 1 份。\hfill {\small\textit{完成标志：交付给队友（对应 Agent 主阶段 0 依赖）}}
\end{itemize}

\subsection*{阶段 1 · 领域知识 + 造题（第 2-3 周）}

\begin{itemize}
\item[☐] 按 schema 整理常见 HLS 失败模式与修法（先覆盖编译错+csim 功能 bug）。\hfill {\small\textit{完成标志：≥10 条条目，每条含症状/错误签名/根因/修法/示例}}
\item[☐] 造 3-5 道本地小题：拿官方 example 故意改坏（删 pragma、引入 bug、制造死锁）。\hfill {\small\textit{完成标志：每题有“坏代码+预期修法+已知正确答案”}}
\item[☐] 验证 Agent 改出来的代码确实对：跑 csim/cosim 核对。\hfill {\small\textit{完成标志：能判定 agent 的提交是否 correctness 通过}}
\item[☐] 读 AutoChip / RTLFixer / Automated C/C++ Repair for HLS 三篇。\hfill {\small\textit{完成标志：从中提炼 5 条“修法模式”补充进知识库}}
\end{itemize}

\subsection*{阶段 2 · cosim + 知识库扩充（第 4-5 周）→ 支撑里程碑}

\begin{itemize}
\item[☐] 扩知识库到 cosim 协议类：dataflow 死锁规则、AXI-Stream 握手(TLAST/TREADY/TVALID)、ap\_ctrl 时序。\hfill {\small\textit{完成标志：≥5 条 cosim 类条目}}
\item[☐] 整理 PPA 优化杠杆清单（pipeline/unroll/array\_partition/dataflow/bind\_op）及何时用。\hfill {\small\textit{完成标志：1 份“PPA 优化速查”，供后续用}}
\item[☐] 配合 Agent 主做端到端测试：提供 5-10 道覆盖各类错误的题 + 标准答案。\hfill {\small\textit{完成标志：测试题集交付，参与判定成功率}}
\item[☐] 里程碑汇合：与 Agent 主一起完成“第一个 Agent 跑通”演示。\hfill {\small\textit{完成标志：演示中你的知识库条目至少被命中并生效 1 次}}
\end{itemize}

\section{两人对齐里程碑（汇合点）}

\begin{longtable}{@{}p{0.08\linewidth}p{0.12\linewidth}p{0.26\linewidth}p{0.26\linewidth}p{0.24\linewidth}@{}}
\toprule
\textbf{汇合点} & \textbf{时间} & \textbf{你交付} & \textbf{队友交付} & \textbf{对齐动作} \\
\midrule
\endhead
1 & 第一周末 & API 通路 & 日志样例包 & 你据此设计日志解析器 \\
2 & 第三周末 & 知识库注入接口+评测脚本 & 知识库初版+3-5 道小题 & 跑首次本地评测，看 success rate \\
3 & 第五周末 & 里程碑版 agent & 端到端测试题集 & 一起跑里程碑演示，达成 DoD \\
\bottomrule
\end{longtable}

% ======================================================================
\chapter{风险与待确认事项}

\section{必须尽早确认的比赛规则}

\begin{itemize}
\item 可用模型范围：能否自由用商用 API？还是限定开源模型？这直接决定要不要微调。
\item LLM 调用是否也计预算：若只计 csim/cosim/synth，可放开让强模型多推理；若 LLM 也计额度，策略完全不同。
\item 评估接口的具体形式：怎么调、返回什么报告格式、预算怎么计——拿到后第一时间写 mock 对齐。
\item Track B 的内容：你提到“两条互补赛道”，B 是不是另一面（比如要亲手做 HLS 设计）？若有 B 描述，方向可能要重新权衡。
\end{itemize}

\begin{warnbox}
\textbf{行动项}\par
把以上四条列成清单，主动找比赛方/通知确认。确认结果直接改本手册的战略假设。
\end{warnbox}

\section{技术风险清单}

\begin{longtable}{@{}p{0.20\linewidth}p{0.34\linewidth}p{0.42\linewidth}@{}}
\toprule
\textbf{风险} & \textbf{影响} & \textbf{对策} \\
\midrule
\endhead
HLS 领域知识不足 & agent 瞎猜烧预算 & 队友专攻，沉淀知识库；不指望 LLM 现场悟 \\
Vitis HLS 日志噪声大 & 解析器抽不出关键错误 & 队友阶段 0 先交真实日志样例，据此写解析器 \\
评估接口未发布 & agent 没靶子打 & 先写 mock，真接口一到立刻替换 \\
商用 API 成本 & 烧钱 & 开发期用便宜模型，评测期才上强模型 \\
本地无 Vitis HLS & 无法复现日志/调解析器 & 第一周就装好，别拖 \\
\bottomrule
\end{longtable}

\section{范式争议：先正确 vs 联合优化}

赛道规则要求先修正确性再优化 PPA。但 C-3 的 COEVO 反对这个串行范式，主张联合优化。建议：遵守赛道规则的顺序（这是评分要求），但读 COEVO 让你在设计 agent 的“何时切换到 PPA 阶段”这个判断时不被单一思路锁死——也许 correctness 基本稳住后就可以边优化边回归验证，而不是死等 100\% 正确。

% ======================================================================
\appendix
\chapter{术语表}

\begin{longtable}{@{}p{0.14\linewidth}p{0.40\linewidth}p{0.42\linewidth}@{}}
\toprule
\textbf{术语} & \textbf{展开} & \textbf{一句话} \\
\midrule
\endhead
HLS & High-Level Synthesis & C/C++ → FPGA 硬件 \\
LLM & Large Language Model & agent 的大脑 \\
agent & Autonomous Agent & 能循环调工具完成任务的 LLM 程序 \\
csim & C Simulation & 跑 C 验功能 \\
csynth & C Synthesis & C→RTL 综合并出报告 \\
cosim & C/RTL Co-Simulation & RTL+测试台协同仿真 \\
PPA & Performance/Power/Area & 性能/功耗/资源 \\
II & Initiation Interval & 流水线每隔几拍能输入一个新数据 \\
pragma & HLS Directive & C 里指挥硬件生成的指令 \\
dataflow & Task-Level Pipelining & 函数间任务级流水 \\
RAG & Retrieval-Augmented Generation & 检索知识再生成 \\
DoD & Definition of Done & 完成定义 \\
\bottomrule
\end{longtable}

\chapter{知识库条目 schema 模板}

队友照此填，你照此读。字段含义见第三章 3.3。

\begin{lstlisting}
id:          KB-018
症状关键词:   ["cosim deadlock", "dataflow"]
错误签名:     "deadlock detected.*dataflow"
根因:         "dataflow 流上读写顺序不当/多生产者多消费者"
修法策略:     "检查每条 stream 单生产者单消费者;
              调整读写顺序; 必要时加 ping-pong 缓冲"
代码示例:     "#pragma HLS dataflow
              #pragma HLS stream depth=2 variable=s"
涉及pragma:   ["dataflow", "stream"]
PPA影响:      "修复死锁，可能略增缓冲资源"
阶段:         "correctness"
\end{lstlisting}

\chapter{环境与工具清单}

\begin{longtable}{@{}p{0.16\linewidth}p{0.34\linewidth}p{0.14\linewidth}p{0.32\linewidth}@{}}
\toprule
\textbf{工具} & \textbf{用途} & \textbf{谁装} & \textbf{备注} \\
\midrule
\endhead
Vitis HLS & 跑 csim/csynth/cosim、出报告 & 队友 & AMD 官网下载，需对应 FPGA 平台版本 \\
Node.js + npm & 跑 docx/脚本 & 你 & 已装 \\
Python 3 & 跑 Vitis HLS Tcl 包装、数据处理 & 你+队友 & 本机暂无，需装 \\
Git & 两人共享代码/知识库 & 两人 & 用同一个仓库 \\
商用 LLM API key & agent 大脑 & 你 & GPT/Claude 等 \\
VS Code + 插件 & 写代码 & 两人 & C/C++、Python、Verilog 高亮 \\
\bottomrule
\end{longtable}

\begin{warnbox}
\textbf{收尾}\par
装环境最拖进度。第一周内把 Vitis HLS 和 Python 装好，别让它卡住后面所有事。
\end{warnbox}

\end{document}